\documentclass{article}
\usepackage{graphicx} % Required for inserting images
\usepackage{amsmath}
\usepackage{amsthm}
\usepackage{amssymb}
\usepackage{mathrsfs}
\usepackage{enumitem}
\newtheorem{theorem}{Theorem}
\newtheorem{lemma}{Lemma}
\newtheorem{corollary}{Corollary}
\theoremstyle{remark}\newtheorem{remark}{Remark}

% Theorem with manually specified number, to match Goldrei's numbering
\theoremstyle{plain}
\newtheorem{bookthminner}{Theorem}
\newenvironment{bookthm}[1]
  {\renewcommand\thebookthminner{#1}\bookthminner}
  {\endbookthminner}

\newcommand{\T}[0]{\mathcal{T}}
\newcommand{\B}[0]{\mathcal{B}}
\newcommand{\C}[0]{\mathcal{C}}
\newcommand{\R}[0]{\mathbb{R}}
\newcommand{\Q}[0]{\mathbb{Q}}
\newcommand{\Z}[0]{\mathbb{Z}}
\newcommand{\N}[0]{\mathbb{N}}
\newcommand{\e}[0]{\epsilon}

\begin{document}

2.35 Show each of following, for all formulas $\phi, \psi, \theta$.

(a) $\phi \equiv \phi$

(b) If $\phi \equiv \psi$ then $\psi \equiv \phi$.

(c) If $\phi \equiv \psi$ and $\psi \equiv \theta$, then $\phi \equiv \theta$.

\begin{proof}[Proof of (a)]
    Let $v$ be any truth assignment. Then $v(\phi) = v(\phi)$ so $\phi \equiv \phi$.
\end{proof}

\begin{proof}[Proof of (b)]
    Suppose $\phi \equiv \psi$, by definition we then have $v(\phi) = v(\psi)$ for all truth assignments $v$. So then it is also true that $v(\psi) = v(\phi)$ for all truth assignments $v$. Thus, $\psi \equiv \phi$.
\end{proof}

\begin{proof}[Proof of (c)]
    Suppose $\phi \equiv \psi$ and $\psi \equiv \theta$. By definition, $v(\phi) = v(\psi)$ and $v(\psi) = v(\theta)$ for all truth assignments $v$. So then it is also true that $v(\phi) = v(\theta)$ for all truth assignments $v$ and so $\phi \equiv \theta$.
\end{proof}

2.34 Show that $\phi \equiv \psi$ if and only if $(\phi \iff \psi)$ is a tautology, for all formulas $\phi, \psi$.

\begin{proof}
    ($\Rightarrow$):
    Suppose $\phi \equiv \psi$. By definition for all truth assignments, $v$, $v(\phi) = v(\psi)$. Let $v$ be any truth assignment. Then $v(\phi) = v(\psi)$, so by the truth table for $\iff$ $v((\phi \iff \psi)) = T$. Since $v$ was arbitrary, $(\phi \iff \psi)$ is a tautology.

    ($\Leftarrow$):
    Suppose $(\phi \iff \psi)$ is a tautology. Then for any truth assignment $v$, we have $v((\phi \iff \psi)) = T$. The truth table shows $\iff$ takes value $T$ only when both sides agree. Hence $v(\phi) = v(\psi)$. Since $v$ was arbitrary, $\phi \equiv \psi$.
\end{proof}

\end{document}
